% Options for packages loaded elsewhere
\PassOptionsToPackage{unicode}{hyperref}
\PassOptionsToPackage{hyphens}{url}
\documentclass[
]{article}
\usepackage{xcolor}
\usepackage[margin=1in]{geometry}
\usepackage{amsmath,amssymb}
\setcounter{secnumdepth}{-\maxdimen} % remove section numbering
\usepackage{iftex}
\ifPDFTeX
  \usepackage[T1]{fontenc}
  \usepackage[utf8]{inputenc}
  \usepackage{textcomp} % provide euro and other symbols
\else % if luatex or xetex
  \usepackage{unicode-math} % this also loads fontspec
  \defaultfontfeatures{Scale=MatchLowercase}
  \defaultfontfeatures[\rmfamily]{Ligatures=TeX,Scale=1}
\fi
\usepackage{lmodern}
\ifPDFTeX\else
  % xetex/luatex font selection
\fi
% Use upquote if available, for straight quotes in verbatim environments
\IfFileExists{upquote.sty}{\usepackage{upquote}}{}
\IfFileExists{microtype.sty}{% use microtype if available
  \usepackage[]{microtype}
  \UseMicrotypeSet[protrusion]{basicmath} % disable protrusion for tt fonts
}{}
\makeatletter
\@ifundefined{KOMAClassName}{% if non-KOMA class
  \IfFileExists{parskip.sty}{%
    \usepackage{parskip}
  }{% else
    \setlength{\parindent}{0pt}
    \setlength{\parskip}{6pt plus 2pt minus 1pt}}
}{% if KOMA class
  \KOMAoptions{parskip=half}}
\makeatother
\usepackage{color}
\usepackage{fancyvrb}
\newcommand{\VerbBar}{|}
\newcommand{\VERB}{\Verb[commandchars=\\\{\}]}
\DefineVerbatimEnvironment{Highlighting}{Verbatim}{commandchars=\\\{\}}
% Add ',fontsize=\small' for more characters per line
\usepackage{framed}
\definecolor{shadecolor}{RGB}{248,248,248}
\newenvironment{Shaded}{\begin{snugshade}}{\end{snugshade}}
\newcommand{\AlertTok}[1]{\textcolor[rgb]{0.94,0.16,0.16}{#1}}
\newcommand{\AnnotationTok}[1]{\textcolor[rgb]{0.56,0.35,0.01}{\textbf{\textit{#1}}}}
\newcommand{\AttributeTok}[1]{\textcolor[rgb]{0.13,0.29,0.53}{#1}}
\newcommand{\BaseNTok}[1]{\textcolor[rgb]{0.00,0.00,0.81}{#1}}
\newcommand{\BuiltInTok}[1]{#1}
\newcommand{\CharTok}[1]{\textcolor[rgb]{0.31,0.60,0.02}{#1}}
\newcommand{\CommentTok}[1]{\textcolor[rgb]{0.56,0.35,0.01}{\textit{#1}}}
\newcommand{\CommentVarTok}[1]{\textcolor[rgb]{0.56,0.35,0.01}{\textbf{\textit{#1}}}}
\newcommand{\ConstantTok}[1]{\textcolor[rgb]{0.56,0.35,0.01}{#1}}
\newcommand{\ControlFlowTok}[1]{\textcolor[rgb]{0.13,0.29,0.53}{\textbf{#1}}}
\newcommand{\DataTypeTok}[1]{\textcolor[rgb]{0.13,0.29,0.53}{#1}}
\newcommand{\DecValTok}[1]{\textcolor[rgb]{0.00,0.00,0.81}{#1}}
\newcommand{\DocumentationTok}[1]{\textcolor[rgb]{0.56,0.35,0.01}{\textbf{\textit{#1}}}}
\newcommand{\ErrorTok}[1]{\textcolor[rgb]{0.64,0.00,0.00}{\textbf{#1}}}
\newcommand{\ExtensionTok}[1]{#1}
\newcommand{\FloatTok}[1]{\textcolor[rgb]{0.00,0.00,0.81}{#1}}
\newcommand{\FunctionTok}[1]{\textcolor[rgb]{0.13,0.29,0.53}{\textbf{#1}}}
\newcommand{\ImportTok}[1]{#1}
\newcommand{\InformationTok}[1]{\textcolor[rgb]{0.56,0.35,0.01}{\textbf{\textit{#1}}}}
\newcommand{\KeywordTok}[1]{\textcolor[rgb]{0.13,0.29,0.53}{\textbf{#1}}}
\newcommand{\NormalTok}[1]{#1}
\newcommand{\OperatorTok}[1]{\textcolor[rgb]{0.81,0.36,0.00}{\textbf{#1}}}
\newcommand{\OtherTok}[1]{\textcolor[rgb]{0.56,0.35,0.01}{#1}}
\newcommand{\PreprocessorTok}[1]{\textcolor[rgb]{0.56,0.35,0.01}{\textit{#1}}}
\newcommand{\RegionMarkerTok}[1]{#1}
\newcommand{\SpecialCharTok}[1]{\textcolor[rgb]{0.81,0.36,0.00}{\textbf{#1}}}
\newcommand{\SpecialStringTok}[1]{\textcolor[rgb]{0.31,0.60,0.02}{#1}}
\newcommand{\StringTok}[1]{\textcolor[rgb]{0.31,0.60,0.02}{#1}}
\newcommand{\VariableTok}[1]{\textcolor[rgb]{0.00,0.00,0.00}{#1}}
\newcommand{\VerbatimStringTok}[1]{\textcolor[rgb]{0.31,0.60,0.02}{#1}}
\newcommand{\WarningTok}[1]{\textcolor[rgb]{0.56,0.35,0.01}{\textbf{\textit{#1}}}}
\usepackage{graphicx}
\makeatletter
\newsavebox\pandoc@box
\newcommand*\pandocbounded[1]{% scales image to fit in text height/width
  \sbox\pandoc@box{#1}%
  \Gscale@div\@tempa{\textheight}{\dimexpr\ht\pandoc@box+\dp\pandoc@box\relax}%
  \Gscale@div\@tempb{\linewidth}{\wd\pandoc@box}%
  \ifdim\@tempb\p@<\@tempa\p@\let\@tempa\@tempb\fi% select the smaller of both
  \ifdim\@tempa\p@<\p@\scalebox{\@tempa}{\usebox\pandoc@box}%
  \else\usebox{\pandoc@box}%
  \fi%
}
% Set default figure placement to htbp
\def\fps@figure{htbp}
\makeatother
\setlength{\emergencystretch}{3em} % prevent overfull lines
\providecommand{\tightlist}{%
  \setlength{\itemsep}{0pt}\setlength{\parskip}{0pt}}
\usepackage{bookmark}
\IfFileExists{xurl.sty}{\usepackage{xurl}}{} % add URL line breaks if available
\urlstyle{same}
\hypersetup{
  pdftitle={1.1 Exploratory Data Analysis (EDA)},
  pdfauthor={Thoa Le},
  hidelinks,
  pdfcreator={LaTeX via pandoc}}

\title{1.1 Exploratory Data Analysis (EDA)}
\author{Thoa Le}
\date{}

\begin{document}
\maketitle

\begin{Shaded}
\begin{Highlighting}[]
\FunctionTok{library}\NormalTok{(fpp3) }\CommentTok{\# time series analysis package}
\end{Highlighting}
\end{Shaded}

\begin{verbatim}
## -- Attaching packages -------------------------------------------- fpp3 1.0.3 --
\end{verbatim}

\begin{verbatim}
## v tibble      3.3.0     v tsibble     1.2.0
## v dplyr       1.1.4     v tsibbledata 0.4.1
## v tidyr       1.3.1     v ggtime      0.2.0
## v lubridate   1.9.5     v feasts      0.5.0
## v ggplot2     4.0.0     v fable       0.5.0
\end{verbatim}

\begin{verbatim}
## -- Conflicts ------------------------------------------------- fpp3_conflicts --
## x lubridate::date()    masks base::date()
## x dplyr::filter()      masks stats::filter()
## x tsibble::intersect() masks base::intersect()
## x tsibble::interval()  masks lubridate::interval()
## x dplyr::lag()         masks stats::lag()
## x tsibble::setdiff()   masks base::setdiff()
## x tsibble::union()     masks base::union()
\end{verbatim}

\begin{Shaded}
\begin{Highlighting}[]
\FunctionTok{library}\NormalTok{(tsibble) }\CommentTok{\# tsibble objects}
\FunctionTok{library}\NormalTok{(lubridate) }\CommentTok{\# for data to work with dates and times}
\FunctionTok{library}\NormalTok{(feasts)}
\CommentTok{\#library(latex2exp)}

\CommentTok{\# read in the data}
\NormalTok{data }\OtherTok{\textless{}{-}}\NormalTok{ readr}\SpecialCharTok{::}\FunctionTok{read\_csv}\NormalTok{(}\StringTok{"BirthsAndFertilityRatesAnnual.csv"}\NormalTok{)}
\end{Highlighting}
\end{Shaded}

\begin{verbatim}
## Rows: 17 Columns: 67
\end{verbatim}

\begin{verbatim}
## -- Column specification --------------------------------------------------------
## Delimiter: ","
## chr (22): DataSeries, 2025, 1979, 1978, 1977, 1976, 1975, 1974, 1973, 1972, ...
## dbl (45): 2024, 2023, 2022, 2021, 2020, 2019, 2018, 2017, 2016, 2015, 2014, ...
## 
## i Use `spec()` to retrieve the full column specification for this data.
## i Specify the column types or set `show_col_types = FALSE` to quiet this message.
\end{verbatim}

\section{obtaining the data
specified}\label{obtaining-the-data-specified}

BirthAndFertility \textless- data \textbar\textgreater{} pivot\_longer(
\# change rows to column cols = \texttt{1960}:\texttt{2025}, names\_to =
``Year'', values\_to = ``Value'', \# naming it to remember that it needs
to be values values\_transform = list(Value = as.numeric) )
\textbar\textgreater{} filter(DataSeries \%in\% c(``Total Fertility Rate
(TFR)'', ``Total Live-Births'')) \textbar\textgreater{} mutate(Year =
as.integer(Year)) \textbar\textgreater{} filter(Year != 2025)
\textbar\textgreater{} \# assignment specified until 2024 pivot\_wider(
\# taking TFR and Total Live-Births into 2 separate columns according to
year names\_from = DataSeries, values\_from = Value )

\section{created a tsibble}\label{created-a-tsibble}

BirthAndFertility \textless- as\_tsibble(BirthAndFertility, index =
Year)

\section{checking it is a tsibble}\label{checking-it-is-a-tsibble}

ts\_BirthAndFertility is\_tsibble(BirthAndFertility)

\section{Preliminary (exploratory analysis) \{then choosing and fitting
models, then using \& evaluating a forecasting
model\}}\label{preliminary-exploratory-analysis-then-choosing-and-fitting-models-then-using-evaluating-a-forecasting-model}

\section{plot the Total Fertility Rate (TFR) and Total Live-Births
(TLB)}\label{plot-the-total-fertility-rate-tfr-and-total-live-births-tlb}

autoplot(BirthAndFertility, \texttt{Total\ Fertility\ Rate\ (TFR)}) +
labs(title = ``Initial Plot of Total Fertility Rate (TFR)'') \# obvious
decreasing trend \# cannot determine seasonality from yearly data \#
there might be cyclic but will have to take away trend to find out

autoplot(BirthAndFertility,\texttt{Total\ Live-Births})+ labs(title =
``Initial Plot of Total Live-Births'') \# there is a decreasing trend
from 1960 to 2024 \# cannot determine seasonality from yearly data \#
there might also be cyclic but will have to take away trend to find out

\section{evidence of the presence of business
cycles?}\label{evidence-of-the-presence-of-business-cycles}

\section{outliers that need to be explained by
experts?}\label{outliers-that-need-to-be-explained-by-experts}

\section{strength of relationships among
variables?}\label{strength-of-relationships-among-variables}

BirthAndFertility \textbar\textgreater{} ggplot(aes(x =
\texttt{Total\ Fertility\ Rate\ (TFR)}, y =
\texttt{Total\ Live-Births})) + geom\_point() + geom\_smooth() +
labs(title = ``Total Fertility Rate (TFR) vs Total Live-Births (TLB)'')
\# there seems to be a nonlinear positive relationship between TFR and
TLB but need to do further analysis \# noting that higher fertility
rates will increase births

\section{lag plots}\label{lag-plots}

BirthAndFertility \textbar\textgreater{}
gg\_lag(\texttt{Total\ Fertility\ Rate\ (TFR)}, geom = `point') +
labs(title = ``Lag Plot of Total Fertility Rate (TFR)'') \# stronger
positive relationships at lags 1 and 2

BirthAndFertility \textbar\textgreater{}
gg\_lag(\texttt{Total\ Live-Births}, geom = `point') + labs(title =
``Lag Plot of Total Live-Births (TLB)'') \# stronger positive
relationships at lags 1 and 2 as well

\section{autocorrelation}\label{autocorrelation}

BirthAndFertility \textbar\textgreater{}
ACF(\texttt{Total\ Fertility\ Rate\ (TFR)}, lag\_max = 9)
\textbar\textgreater{} autoplot() + labs(title = ``ACF Plot for Total
Fertility Rate (TFR)'') \# lags 1-9 are significant so more lags should
be checked \# slow decrease in the ACF as the lags increase indicates a
trend

BirthAndFertility \textbar\textgreater{}
ACF(\texttt{Total\ Live-Births}, lag\_max = 9) autoplot() + labs(title =
``ACF Plot for Total Live-Births (TLB)'') \# lags 1-6 are significant
suggests short-term dependence \# slow decrease in the ACF as the lags
increase indicates a trend \# as there are more than one large spike
outside the dotted bounds, this series is probably not white noise

BirthAndFertility \textbar\textgreater{}
ACF(\texttt{Total\ Fertility\ Rate\ (TFR)}, lag\_max = 24)
\textbar\textgreater{} autoplot() + labs(title = ``ACF Plot for Total
Fertility Rate (TFR)'') \# lags 1-10 are significant suggests strong
persistence \# slow decrease in the ACF as the lags increase indicates a
trend \# as there are more than one large spike outside the dotted
bounds, this series is probably not white noise

\section{Box-Cox transformation}\label{box-cox-transformation}

lambdatfr \textless- BirthAndFertility \textbar\textgreater{}
features(\texttt{Total\ Fertility\ Rate\ (TFR)}, features = guerrero)
\textbar\textgreater{} \# Guerrero method pull(lambda\_guerrero) \#
Box-Cox transformation parameter = -0.44514 \# as lambda is below zero
but halfway to -1 suggests moderate transformation required
BirthAndFertility \textbar\textgreater{}
autoplot(box\_cox(\texttt{Total\ Fertility\ Rate\ (TFR)}, lambdatfr)) +
labs(title = latex2exp::TeX(paste0( ``Transformed Total Fertility Rate
(TFR) with \(\\lambda\) ='', round(lambdatfr,2)))) \# decreasing trend
still obvious \# still signs there might also be cyclic but still will
have to take away trend to find out

lambdatlb \textless- BirthAndFertility \textbar\textgreater{}
features(\texttt{Total\ Live-Births}, features = guerrero)
\textbar\textgreater{} \# Guerrero method pull(lambda\_guerrero) \#
Box-Cox transformation parameter = -0.89992 \# as lambda is below zero
but close to -1 suggests strong transformation required
BirthAndFertility \textbar\textgreater{}
autoplot(box\_cox(\texttt{Total\ Live-Births}, lambdatlb)) + labs(title
= latex2exp::TeX(paste0( ``Transformed Total Live-Births (TLB) with
\(\\lambda\) ='', round(lambdatlb,2)))) \# the decreasing trend from
1960 to 2024 is less prominent \# might still be cyclic but still will
have to take away trend to find out

\section{STL decomposition method:
trend}\label{stl-decomposition-method-trend}

dcmptfr \textless- BirthAndFertility \textbar\textgreater{} model(stl =
STL(\texttt{Total\ Fertility\ Rate\ (TFR)})) components(dcmptfr)
\textbar\textgreater{} as\_tsibble() \textbar\textgreater{}
autoplot(\texttt{Total\ Fertility\ Rate\ (TFR)}), colour = ``grey'') +
\# raw data geom\_line(aes(y = trend), colour = ``\#D55E00'') + \#
trend-cycle component labs(title = ``Plot of trend-cycle component
(orange) and raw data (grey) for TFR'') \# little variability to the raw
data, more smooth

dcmptlb \textless- BirthAndFertility \textbar\textgreater{} model(stl =
STL(\texttt{Total\ Live-Births})) components(dcmptlb)
\textbar\textgreater{} as\_tsibble() \textbar\textgreater{}
autoplot(\texttt{Total\ Live-Births}), colour = ``grey'') + \# raw data
geom\_line(aes(y = trend), colour = ``\#D55E00'') + \# trend-cycle
component labs(title = ``Plot of trend-cycle component (orange) and raw
data (grey) for TLB'') \# more variability to the raw data, bigger
changes

\section{STL decomposition}\label{stl-decomposition}

components(dcmptfr) \textbar\textgreater{} autoplot() \# decreasing
trend as mentioned earlier

components(dcmptlb) \textbar\textgreater{} autoplot() \# overall
decreasing trend but there was an increase from 1980 with a peak in 1991
before decreasing again

\section{the remainder of both decompositions resemble each other with
sharp dip around 1986 and spike
1988}\label{the-remainder-of-both-decompositions-resemble-each-other-with-sharp-dip-around-1986-and-spike-1988}

\section{moving average}\label{moving-average}

BirthAndFertility\_MA\_TFR \textless- BirthAndFertility
\textbar\textgreater{} mutate(\texttt{5-MA\_TFR} =
slider::slide\_dbl(\texttt{Total\ Fertility\ Rate\ (TFR)}, mean, .before
= 2, .after = 2, .complete = TRUE)) BirthAndFertility\_MA\_TFR
\textbar\textgreater{} autoplot(\texttt{Total\ Fertility\ Rate\ (TFR)})
+ geom\_line(aes(y = \texttt{5-MA\_TFR}), colour = ``\#D55E00'') +
labs(title = ``Total Fertility Rate (black) and 5-MA estimate of
trend-cycle (orange)'') \# the original data was already quite smooth so
it matches quite closely with the trend-cycle

BirthAndFertility\_MA\_TLB \textless- BirthAndFertility
\textbar\textgreater{} mutate(\texttt{5-MA\_TLB} =
slider::slide\_dbl(\texttt{Total\ Live-Births}, mean, .before = 2,
.after = 2, .complete = TRUE)) BirthAndFertility\_MA\_TLB
\textbar\textgreater{} autoplot(\texttt{Total\ Live-Births}) +
geom\_line(aes(y = \texttt{5-MA\_TLB}), colour = ``\#D55E00'') +
labs(title = ``Total Live-Births (black) and 5-MA estimate of
trend-cycle (orange)'') \# can see smooth curve but due to the roughness
of the data should increase the average for a smoother curve

BirthAndFertility\_MA\_TLB7 \textless- BirthAndFertility
\textbar\textgreater{} mutate(\texttt{7-MA\_TLB} =
slider::slide\_dbl(\texttt{Total\ Live-Births}, mean, .before = 3,
.after = 3, .complete = TRUE)) BirthAndFertility\_MA\_TLB7
\textbar\textgreater{} autoplot(\texttt{Total\ Live-Births}) +
geom\_line(aes(y = \texttt{7-MA\_TLB}), colour = ``\#D55E00'') +
labs(title = ``Total Live-Births (black) and 7-MA estimate of
trend-cycle (orange)'') \# the curve is smoother at the cost of the ends

\section{cannot do classical decomposition as data is yearly and need at
least 2
periods}\label{cannot-do-classical-decomposition-as-data-is-yearly-and-need-at-least-2-periods}

\section{BirthAndFertility
\textbar\textgreater{}}\label{birthandfertility}

\section{\texorpdfstring{model(classical\_decomposition(\texttt{Total\ Fertility\ Rate\ (TFR)},
type =
``additive''))}{model(classical\_decomposition(Total Fertility Rate (TFR), type = ``additive''))}}\label{modelclassical_decompositiontotal-fertility-rate-tfr-type-additive}

\section{components() \textbar\textgreater{}}\label{components}

\section{\texorpdfstring{autoplot(\texttt{Total\ Fertility\ Rate\ (TFR)})
+}{autoplot(Total Fertility Rate (TFR)) +}}\label{autoplottotal-fertility-rate-tfr}

\section{labs(title = ``Classical additive decomposition of Total
Fertility
Rate'')}\label{labstitle-classical-additive-decomposition-of-total-fertility-rate}

\section{also cannot use the X-11 method for the same reason as
classical
decomposition}\label{also-cannot-use-the-x-11-method-for-the-same-reason-as-classical-decomposition}

\end{document}
